The minimum subsidy is $S^*_m=\max\{\hat R-\bar R_m,0\}$, evaluated at two renovation-cost benchmarks: \euro{}250/m\textsuperscript{2} (a light, ``geste par geste'' intervention such as roof or wall insulation) and \euro{}500/m\textsuperscript{2} (a medium, two-measure upgrade), drawing on MaPrimeR\'enov' programme documentation.\footnote{MaPrimeR\'enov' distinguishes ``geste par geste'' interventions (insulation of walls or roof, costing roughly \euro{}100--400/m\textsuperscript{2} of treated area) from comprehensive ``r\'enovation d'ampleur'' operations (\euro{}700--1{,}200/m\textsuperscript{2} or more); the \euro{}500/m\textsuperscript{2} benchmark corresponds broadly to a two-measure upgrade.} The benchmarks are illustrative. The light benchmark corresponds to a partial upgrade and not to the full move from an E--G to an A--D rating that the premium $b$ prices, and the medium benchmark only approaches it; a comprehensive renovation at the cost range of the footnote would exceed the threshold in every market. A single $\hat R$ per market predicts that no informed low-quality seller exists where $\hat R<\hat b$, yet low ratings are disclosed in every market, so effective renovation costs exceed the benchmark in the \attBaseNoSubsidyTwoFiftyApt{} percent of apartment and \attBaseNoSubsidyTwoFiftyHouse{} percent of house markets where it falls below the premium. At $\hat R=250$, \attScBaseShareTwoFiftyApt{} percent of apartment markets require a positive subsidy; conditional on requiring one, the median $S^*_m$ is \euro{}\attBaseSMedTwoFiftyApt{}/m\textsuperscript{2}, or about \euro{}\attBaseUnitTwoFiftyApt{} per unit at the median apartment area of \attBaseAreaMedApt{}~m\textsuperscript{2}. For houses, \attScBaseShareTwoFiftyHouse{} percent of markets require support, with a conditional median of \euro{}\attBaseSMedTwoFiftyHouse{}/m\textsuperscript{2} (about \euro{}\attBaseUnitTwoFiftyHouse{} per unit at \attBaseAreaMedHouse{}~m\textsuperscript{2}). At $\hat R=500$ the shares are \attScBaseShareFiveHundredApt{} and \attScBaseShareFiveHundredHouse{} percent and the conditional medians \euro{}\attBaseSMedFiveHundredApt{} and \euro{}\attBaseSMedFiveHundredHouse{}/m\textsuperscript{2}. The cost-sharing ratio $S^*_m/\hat R$ is the natural decision variable for the regulator. Among markets requiring support its mean is \attBaseCostShareTwoFiftyApt{} percent for apartments and \attBaseCostShareTwoFiftyHouse{} percent for houses at the light benchmark, and \attBaseCostShareFiveHundredApt{} and \attBaseCostShareFiveHundredHouse{} percent at the medium one. A commitment to a fixed rate is therefore more defensible than a commitment to a fixed euro amount, because it adjusts to the depth of the intervention and to local conditions.
